\chapter{How Belief Regimes Fracture}
\label{ch:how-belief-regimes-fracture}

\epigraph{There are decades where nothing happens; and there are weeks where decades happen.}{Vladimir Lenin (attr.)}

\noindent
Doxonomic equilibria are stable but not permanent.
This chapter analyzes the mechanisms by which dominant belief regimes fracture: contradictions, scandals, economic shocks, institutional failure, and the loss of elite credibility.

\section{Regime Fracture as Bifurcation}
\label{sec:ch30-bifurcation}

\begin{model}[Bifurcation Model of Regime Fracture]
\label{mod:bifurcation}
\index{bifurcation}
Let $\belief{}(t)$ be the strength of the dominant belief and $\sigma(t)$ a stress parameter (accumulating contradictions, declining legitimacy).
The dynamics are:
\[
  \frac{d\belief{}}{dt} = \belief{}(\alpha\phi - \delta - \sigma(t)) + \beta\belief{}^2(1 - \belief{})
\]
As $\sigma$ increases past a critical value $\sigma^*$, the high-$\belief{}$ equilibrium disappears through a saddle-node bifurcation.
The system drops to a low-$\belief{}$ state---regime fracture.
\end{model}

\begin{definition}[Credibility Collapse]
\label{def:credibility-collapse}
\index{credibility collapse}
A \emph{credibility collapse} occurs when the credibility $\credibility{k}$ of a central institution drops rapidly, typically due to a scandal, failure, or exposure of deception.
Because credibility compounds through the network, a single collapse can cascade.
\end{definition}

\section{Sources of Stress}
\label{sec:ch30-stress}

\begin{enumerate}[nosep]
  \item \textbf{Contradictions}: the belief system predicts $X$ but reality delivers $\neg X$ (e.g., ``markets self-correct'' during a crash).
  \item \textbf{Scandals}: institutional hypocrisy is exposed (e.g., a think tank's hidden funding revealed).
  \item \textbf{Economic shocks}: material conditions deteriorate, undermining narrative plausibility.
  \item \textbf{Memetic overload}: the narrative space becomes saturated with competing frames, raising entropy.
  \item \textbf{Generational turnover}: new cohorts arrive without the same formative exposure.
\end{enumerate}

\section{Notes and References}
\label{sec:ch30-notes}

On the dynamics of regime change, see \textcite{gramsci1971} on interregnum.
Threshold and cascade dynamics are in \textcite{granovetter1978} and \textcite{watts2002}.
On economic shocks and ideological change, see \textcite{polanyi1944} and \textcite{klein2007}.

\section{Exercises}

\begin{exercise}
For the bifurcation model, compute $\sigma^*$ as a function of $\alpha$, $\phi$, $\delta$, and $\beta$.
Interpret: what makes a regime more or less resilient to stress?
\end{exercise}

\begin{exercise}
Analyze the 2008 financial crisis as a partial fracture of the ``efficient markets'' belief regime.
Why was the fracture incomplete?
\end{exercise}

\begin{exercise}
Model a credibility cascade in which the exposure of one institution's deception reduces credibility of connected institutions.
Use network contagion dynamics.
\end{exercise}
